The separation also applies to arbitrary curve potentials and smooth obstacles that split by coordinates. Keep the curves $C_i$, the product $Y$ and its projections $\pi_i$. Let $\sigma_i$ be smooth real $(1,1)$-forms with $\PSH(C_i,\sigma_i)\ne\varnothing$, and choose $u_i\in\PSH(C_i,\sigma_i)$ and $F_i\in C^\infty(C_i,\mathbb R)$. Put
\[
 \gamma=\sum_i\pi_i^*\sigma_i,\qquad
 u=\sum_i u_i\circ\pi_i,\qquad
 F=\sum_i F_i\circ\pi_i.
\]

\Needspace{14\baselineskip}
\begin{proposition}\label{prop:separated-curve-obstacle-projections}\index{Envelope!separated smooth obstacles}
For the preceding data,
\begin{equation}\label{eq:separated-curve-obstacle-projections}
 \begin{aligned}
 P_{\gamma,F}[u]
     &=\sum_i P_{\sigma_i,F_i}[u_i]\circ\pi_i,\\
 P_{\gamma,F}[u]_{\mathcal I}
     &=\sum_i P_{\sigma_i,F_i}[u_i]_{\mathcal I}\circ\pi_i.
 \end{aligned}
\end{equation}
These identities include curve backgrounds of zero area.
\end{proposition}
\begin{proof}
Write $e_i=P_{\sigma_i,F_i}[u_i]$. Sums of coordinatewise ordinary competitors are product competitors. Their regularized supremum is $\sum_i e_i\circ\pi_i$, which gives the lower bound for the first identity.

For the upper bound use induction on the number of factors, with the one-factor identity as the initial case. Let $v\leq F$ be an ordinary competitor, so $v\leq u+C$ for some constant $C$. Choose $C_0$ with $u-C_0\leq F$ and replace $v$ by $\max(v,u-C_0)$. This preserves its bounds and makes the required slices genuine.

Write $Y=Y'\times C_m$, and on $Y'$ put $u'=\sum_{i<m}u_i\circ\pi_i$ and $F'=\sum_{i<m}F_i\circ\pi_i$. The regularized partial supremum
\[
 a(t)=\left(\sup_{y'\in Y'}\bigl(v(y',t)-F'(y')\bigr)\right)^*
\]
is $\sigma_m$-plurisubharmonic by the local upper-envelope theorem of \cref{chap:framework}. Indeed, its defining restrictions have the common bound $F_m$, and one is genuine by the lower bound $v\geq u-C_0$. The bound $v\leq u+C$ gives
\[
 a\leq u_m+C+\sup_{Y'}(u'-F'),\qquad a\leq F_m.
\]
Thus $a\leq e_m$. If $u_m(t)$ is finite, so is $e_m(t)$: a sufficiently low translate of $u_m$ is a coordinate competitor. The function $v(\cdot,t)-e_m(t)$ is then a genuine potential in the product background on $Y'$, below $F'$ and more singular than $u'$. Induction gives $v(\cdot,t)\leq\sum_{i<m}e_i\circ\pi_i+e_m(t)$. At a pole of $u_m$, the original bound $v\leq u+C$ already gives $v=-\infty$. Every competitor is therefore bounded by $\sum_i e_i\circ\pi_i$. Taking the regularized supremum proves the ordinary identity.

For the ideal identity put $B_i=\int_{C_i}\sigma_i\geq0$ and $\widehat\eta_i=\eta_i/A_i$. By \cref{lem:curve-background-area}, choose smooth $\chi_i$ with
\[
 \sigma_i=B_i\widehat\eta_i+\ddc\chi_i.
\]
Let $E_i$ be the atomic part of $\sigma_i+\ddc u_i$. If $B_i>0$, apply \cref{lem:curve-atomic-ideal-model} in the positive background $B_i\widehat\eta_i$ to $u_i+\chi_i$. Subtracting $\chi_i$ from its atomic Green potential gives a $\sigma_i$-plurisubharmonic function $g_i$ such that
\[
 \begin{gathered}
 u_i\sim_{\mathcal I}g_i,\qquad
 P_{\sigma_i}[u_i]_{\mathcal I}\sim g_i,\\
 \ddc g_i=E_i-E_i(C_i)\widehat\eta_i-\ddc\chi_i.
 \end{gathered}
\]
The last equivalence uses the smooth-gauge bounds of \cref{prop:i-envelope-gauge}. If $B_i=0$, the area lemma gives $u_i=-\chi_i+c_i$ and $E_i=0$; taking $g_i=-\chi_i$ gives the same conclusions. In either case
\[
 \sigma_i+\ddc g_i
       =E_i+(B_i-E_i(C_i))\widehat\eta_i\geq0.
\]

Put $g=\sum_i g_i\circ\pi_i$. Pullback and addition in \cref{lem:submersion-multiplier-pullback,lem:hermitian-i-calculus} give $g\sim_{\mathcal I}u$. Its Hessian is a countable sum of coordinate-fiber divisor currents minus a smooth form. This divisor current has finite mass: each coordinate fiber has constant mass for a fixed product metric, and each $E_i$ has finite mass. Thus \cref{lem:divisorial-potential-domination} applies even when some $B_i$ vanish. Every ideal competitor $v\leq F$ for $u$ satisfies $v\preceq g$, and is bounded by $P_{\gamma,F}[g]$.

Set $r_i=P_{\sigma_i,F_i}[u_i]_{\mathcal I}$. Ordinary-type invariance and \cref{prop:ideal-obstacle-reprojection} give
\[
 P_{\sigma_i,F_i}[g_i]
    =P_{\sigma_i,F_i}[P_{\sigma_i}[u_i]_{\mathcal I}]
    =r_i.
\]
The proved ordinary separation therefore bounds the ideal envelope above by $\sum_i r_i\circ\pi_i$. Conversely, each $r_i$ is below $F_i$ and ideally equivalent to $u_i$, again by \cref{prop:ideal-obstacle-reprojection}. Their sum is a $\gamma$-plurisubharmonic ideal competitor for $u$ below $F$, proving the reverse inequality. This proves the ideal identity with no strict positivity needed in the divisor comparison.
\end{proof}

The obstacle in \eqref{eq:separated-curve-obstacle-projections} is required to split into curve functions. An arbitrary smooth function on $Y$ does not supply the coordinatewise bounds used in the proof.
